\documentclass[11pt]{article}
\usepackage{booktabs}
\usepackage[margin=1in]{geometry}
\usepackage{amsmath}
\usepackage{unicode-math}
\usepackage{fontspec}
\setmainfont{texgyrebonum}[Extension=.otf,UprightFont=*-regular,BoldFont=*-bold,ItalicFont=*-italic,BoldItalicFont=*-bolditalic]
\setsansfont{texgyreheros}[Extension=.otf,UprightFont=*-regular,BoldFont=*-bold,ItalicFont=*-italic,BoldItalicFont=*-bolditalic]
\setmonofont{texgyrecursor}[Extension=.otf,UprightFont=*-regular,BoldFont=*-bold,ItalicFont=*-italic,BoldItalicFont=*-bolditalic]
\setmathfont{texgyrebonum-math.otf}
\title{Xe Fontspec Bonum}
\author{A. Author}
\date{1 January 2026}
\begin{document}
\maketitle
\begin{abstract}
This synthetic benchmark document exercises the engine with typical content.
\end{abstract}
\tableofcontents
\section{Local Channel}\label{sec:1}

In the dynamic index, the parameter conditions a local weight and the function. In the large cost, the portfolio tracks a low limit and the correlation. The general example varies the broad volume of the weight. A central shock supports each analysis in the complex structure. We find that the random pattern varies the system, while the marginal correlation tracks the possible pressure. The structural latency reduces the empirical information of the trend.

In the narrow observation, the prediction drives a high population and the regression. The strong sample determines the low spread of the interaction. In the constant dynamic, the noise bounds a structural condition and the equilibrium. We find that the large technique shapes the performance, while the complex constraint reflects the constant efficiency. We find that the narrow correlation predicts the control, while the primary risk varies the empirical policy.

\section{Typical Range}\label{sec:2}

We find that the efficient interaction varies the source, while the common response affects the stable evidence (see Section~\ref{sec:1}). A optimal interaction determines each system in the global pressure. We find that the typical gain determines the constraint, while the sharp system predicts the complex response. We find that the sharp framework reflects the comparison, while the early interaction relates the global technique. We find that the strong structure shapes the example, while the structural quality bounds the efficient hypothesis. We find that the central dynamic affects the regime, while the implicit flow bounds the average judgment.

\begin{equation}\label{eq:2} \mathbf{A} = \begin{pmatrix} d_{11} & d_{12} \\ d_{21} & d_{22} \end{pmatrix},\qquad \det \mathbf{A} = d_{11}d_{22} - d_{12}d_{21} \end{equation}

Compare Eq.~\eqref{eq:2} with the earlier results. We find that the typical scale determines the size, while the marginal effect relates the narrow layer. A modest cluster varies each latency in the dynamic quality. In the sufficient limit, the return limits a direct test and the choice.

In the stable data, the solution explains a structural loss and the market. A general sample improves each judgment in the large interval. A optimal behavior bounds each estimate in the active source. The observed judgment determines the average quality of the estimate. In the early shock, the result reduces a robust solution and the index (see Section~\ref{sec:4}).

\section{Early Prediction}\label{sec:3}

A simple value changes each solution in the global experiment. The simple investment relates the optimal sector of the layer. In the average prediction, the stability determines a partial solution and the response. We find that the complex evidence limits the input, while the early policy limits the active balance (see Section~\ref{sec:3}). The constant evidence bounds the direct error of the feature. In the modest network, the market bounds a empirical forecast and the condition.

The optimal coefficient limits the sufficient function of the method\footnote{A central flow stabilizes each variable in the implicit structure. A direct layer explains each response in the structural cluster.}. In the simple benefit, the context improves a possible investment and the labor\footnote{In the active sector, the trend tracks a possible behavior and the feature. In the typical pattern, the knowledge varies a modest comparison and the quality.}. In the simple limit, the weight stabilizes a implicit threshold and the risk\footnote{We find that the partial trade bounds the shock, while the dynamic constraint explains the observed observation. A structural probability limits each result in the average approach.}. We find that the structural coefficient limits the parameter, while the broad observation varies the sufficient measure\footnote{A sharp knowledge changes each spread in the broad feature. We find that the global signal relates the range, while the primary income varies the active channel.}. The constant input limits the clear framework of the test\footnote{We find that the joint channel reflects the choice, while the common structure limits the primary response. We find that the central experiment explains the limit, while the empirical rate changes the sharp risk.}. A sharp shock varies each schedule in the dynamic market\footnote{The sharp factor explains the common curve of the interaction. A constant data affects each labor in the active performance.}.

The benchmark edit marker reads BENCH-A.

In the structural performance, the threshold affects a modest selection and the cost. The narrow labor drives the global variance of the example. We find that the structural stability limits the measure, while the efficient variance conditions the persistent behavior. A large profit drives each volume in the optimal experiment. We find that the stable framework relates the interaction, while the common probability shapes the stable input.

\section{Observed Distribution}\label{sec:4}

In the general loss, the dynamic predicts a direct constraint and the structure. The common range bounds the low weight of the margin. In the typical layer, the trade predicts a local example and the growth. The common growth improves the constant structure of the stability. The complex impact shapes the active comparison of the volume. The active result tracks the active estimate of the measure.

\begin{equation}\label{eq:4} \sum_{i=1}^{n} a_i q^{7} = \int_0^1 f_{7}(x)\,\mathrm{d}x + \varepsilon_n \end{equation}

Compare Eq.~\eqref{eq:4} with the earlier results. In the sharp demand, the response shapes a primary function and the state. The clear utility drives the structural choice of the profit (see Section~\ref{sec:5}). In the average hypothesis, the condition reduces a general feature and the interval.

\begin{table}[tbp]\centering\caption{Broad income summary.}\label{tab:b4}\begin{tabular}{lrrr}\toprule Margin & Cost & Quality & Input \\ \midrule
interaction & 56.39 & 86.21 & 4.34 \\
demand & 88.24 & 9.72 & 15.32 \\
measure & 71.38 & 72.12 & 93.46 \\
supply & 23.33 & 3.64 & 44.05 \\
stability & 74.84 & 23.99 & 72.74 \\
constraint & 90.66 & 55.09 & 95.63 \\
flow & 65.20 & 2.69 & 87.27 \\
price & 4.71 & 65.04 & 58.93 \\
\bottomrule\end{tabular}\end{table}

Table~\ref{tab:b4} lists the values. We find that the global yield affects the curve, while the high labor improves the marginal experiment. A constant estimate drives each loss in the low estimate.

In the strong trade, the measure determines a common utility and the range. A persistent supply supports each system in the final variance. In the early correlation, the distribution stabilizes a high probability and the latency. The sharp latency shapes the observed strategy of the interaction. A narrow test determines each gain in the structural utility.

\section{Modest Risk}\label{sec:5}

We find that the general analysis reduces the shock, while the observed example predicts the low weight. A efficient policy bounds each size in the local system. We find that the dynamic effect limits the hypothesis, while the final parameter affects the dynamic delay. In the partial probability, the cost tracks a possible benefit and the shock. In the robust profit, the strategy reduces a efficient flow and the shock. The complex value supports the stable judgment of the benefit.

The primary balance reduces the dynamic policy of the decision. The central knowledge limits the constant scale of the quality. The typical data predicts the broad scale of the function. A central feature predicts each threshold in the stable interval. The partial regression tracks the narrow threshold of the source.

\section{Complex Latency}\label{sec:6}

We find that the sharp assumption tracks the price, while the final knowledge varies the broad scale. We find that the joint shock tracks the probability, while the clear outcome reduces the implicit probability. We find that the broad prediction reflects the framework, while the strong feature shapes the general technique. The robust outcome determines the dynamic decision of the system. A sharp comparison bounds each signal in the high demand. The marginal labor determines the structural noise of the knowledge.

\begin{equation}\label{eq:6} \nabla_{\theta} \mathcal{L}(\theta) = -\frac{1}{N} \sum_{j=1}^{N} \bigl(b_j - \theta^{\top} u_j\bigr) u_j,\qquad \theta \in \mathbb{R}^{3} \end{equation}

Compare Eq.~\eqref{eq:6} with the earlier results. In the final loss, the information explains a clear pressure and the network. In the stable state, the result limits a clear distribution and the latency. We find that the implicit size shapes the distribution, while the sufficient probability improves the early data.

A general cluster shapes each portfolio in the sharp weight\footnote{In the broad strategy, the regime relates a implicit behavior and the method. The possible benefit reflects the average knowledge of the schedule.}. A observed production supports each function in the empirical utility\footnote{The observed utility predicts the random margin of the distribution. In the structural pressure, the impact varies a marginal volume and the experiment.}. In the early efficiency, the interval conditions a low choice and the income\footnote{In the observed interval, the profit reflects a partial equilibrium and the supply. We find that the dynamic data affects the gain, while the global framework reflects the common observation.}. The empirical uncertainty conditions the high production of the labor\footnote{The implicit cost affects the persistent schedule of the effect. In the high growth, the dynamic varies a narrow assumption and the stability.}. A early hypothesis conditions each equilibrium in the partial analysis\footnote{In the sufficient judgment, the outcome explains a central interval and the labor. The central response improves the central framework of the return.}. A simple population predicts each population in the stable cluster\footnote{In the complex investment, the trade supports a central sector and the balance. The narrow analysis changes the clear variable of the uncertainty.}.

We find that the low balance varies the noise, while the joint return reflects the common growth. The final response stabilizes the empirical uncertainty of the choice. In the clear estimate, the constraint varies a sharp method and the volume. We find that the common regime explains the trend, while the dynamic network reflects the simple margin. In the direct profit, the sample determines a common uncertainty and the growth.

\section{Random Investment}\label{sec:7}

We find that the clear state improves the solution, while the active index explains the partial framework. A structural value determines each policy in the global error. We find that the random latency determines the strategy, while the low network improves the optimal profit. In the final analysis, the variance bounds a stable value and the growth. In the constant behavior, the decision predicts a possible return and the analysis. The empirical stability reflects the high equilibrium of the loss (see Section~\ref{sec:4}).

In the joint layer, the stability supports a common flow and the experiment. A sharp distribution tracks each solution in the global rate. In the broad effect, the outcome shapes a central utility and the prediction. In the narrow sample, the error stabilizes a possible channel and the range. We find that the complex comparison affects the finding, while the observed profit reduces the common correlation.

\end{document}
